\chapter{A Parent That Keeps Its Fresh-Sample Responses}
\label{comp:boundary}

A parent can now be constructed that preserves evolution and the fresh uniform-descendant response interface, through any finite number of supplied binary levels. The dynamics is the reciprocal law potential with its Fisher flow: the original smooth identity-channel KL interaction, with an evaluator that couples belief and model. The children keep distinct public laws. The construction works because a particular class of public patterns is invisible to fresh uniform-descendant sampling, and remains invisible under the inherited evolution. The tree is supplied, and nothing here shows the dynamics selecting it.

\section{Four children and three public signs}

Take a general base variable $r$ with reference law $A$, and three independent public angles. Each angle determines a balanced sign, for example
\begin{equation}
 \sigma_L=\operatorname{sign}\cos(\theta_L-\alpha_L(r)),
\end{equation}
and similarly $\sigma_R,\sigma_0$ for the right node and the root. The phase $\alpha_L(r)$ can be any measurable function. Rotating it changes which actual public samples carry positive or negative contrast. A fair coin with a base-dependent orientation would serve equally well for every exact statement below; the angles are needed only to rotate a phase continuously.

Choose positive profiles $m_+(r),m_-(r)$ and contrast magnitudes $0<a_+(r),a_-(r)<1$, with $\int(m_++m_-)dA/2=1$. Four child densities are
\begin{equation}
 \begin{aligned}
 q_{L\pm}&=m_{\sigma_0}(1\pm a_{\sigma_0}\sigma_L),\\
 q_{R\pm}&=m_{-\sigma_0}(1\pm a_{-\sigma_0}\sigma_R).
 \end{aligned}
 \label{comp:boundary-four}
\end{equation}
They are laws on the full base-and-angle space. There is a separate copy for the model laws. A fixed joint reference $\rho(r,w)$ couples the belief and model bases and is independent of the angles. Every child still compares its full law with every other child through identity channels.

If a response freshly chooses $L+$ or $L-$ with equal probability and then samples that child, its density is $m_{\sigma_0}$. The sign $\sigma_L$ cancels pointwise. Its phase therefore changes the two individual public laws without changing this pair response. The pair nevertheless retains $a_{\sigma_0}^2$: its magnitude influences the future mixture through the KL forces.

The two first parents have response densities $m_{\sigma_0}$ and $m_{-\sigma_0}$. They are distinct when $m_+\ne m_-$. Choosing one parent uniformly cancels $\sigma_0$ and gives the full-coordinate response density
\begin{equation}
 \bar m(r)=\frac{m_+(r)+m_-(r)}2.
\end{equation}
The root also retains the unordered joint pair
\begin{equation}
 \left\{(m_+(r),a_+(r)^2),(m_-(r),a_-(r)^2)\right\}.
\end{equation}
It can discard the phase of $\sigma_0$ while retaining the magnitudes that determine its future. A density value and its contrast must remain paired.

\section{Why the original dynamics closes}

For a tree of depth $d$, assign a balanced public sign to every internal node. A leaf reads one sign on each step of its path. Write the resulting bit string as $u\in\{-1,1\}^d$, and let its density be $f(r,u)$. This profile is an arbitrary positive bounded function; it is not required to remain a product of independent factors.

Two leaves that separate immediately after their first $k$ common edges have the same first $k$ bits and opposite next bits. Their later signs come from separate branches. Define $P_k$ to keep the shared bits, flip the next bit, and average the remaining bits. The exact original peer cost is
\begin{equation}
 D_k(f)=\int f[\log f-P_k\log f]dA dU_d,
\end{equation}
where $U_d$ is the uniform bit law. The source classes have evolving total weights $b_k$. Differentiating the complete energy gives
\begin{equation}
 h=2\log f+V-\sum_k b_k\left[P_k\log f+\frac{P_kf}{f}\right],
 \qquad \dot f=-f(h-\E_fh).
 \label{comp:boundary-force}
\end{equation}
Here $V(r)=-\int g_0(w)\log\rho(r,w)dB(w)$ carries the other sector's base marginal. The ratio $P_kf/f$ comes from being a source in other agents' costs. Omitting it would give a different flow.

All quantities in \cref{comp:boundary-force} use the profile values and their nested grouping. A branch exchange preserves them. At every physical sample point, the source leaves in each class run through every possible suffix bit string, so each leaf's original force equals the profile force \cref{comp:boundary-force} evaluated at that leaf's own path bits. Nothing is averaged over receivers. Both the original energy and the original Fisher metric carry the factor $2^d$; it cancels in the inherited equation. There is no inserted coarse timescale.

The reason this works is symmetry. The fine energy and the Fisher metric do not change when the public angles are relabeled in any way that keeps the base fixed and preserves the angle laws, or when agents are permuted. The family above is exactly the set of configurations left fixed by those relabelings that are compensated by a tree automorphism, so the flow cannot leave it. The phases, and the order of the two branches at each node, are therefore invisible to the energy, the metric, and the admitted records. This is the classical principle of symmetric criticality applied to a gradient flow. What symmetry does not provide is the explicit cost $D_k$, its incoming term, and the formula for the fresh response.

The parent therefore keeps an unordered tree of profile values. Joining siblings places their two retained trees inside one unordered pair. Subsequent joins can be evaluated in the same way, and the result agrees with retaining the entire subtree directly. The grouping must remain: a flat multiset would forget which profiles enter each $D_k$. The parent state can still grow with the number of children.

This joining changes the description of one already coupled fine system. It does not switch on new edges or prove that the agents have selected that tree. Original uniform source rows are allowed; their class totals evolve under the original row equations.

\section{What the inherited interface permits}

At every level, choose a descendant freshly and uniformly, draw its full original sample, and conceal its label. The retained subtree mean gives the exact response density. Choosing a parent and then a descendant gives the same response as direct descendant selection. Repeat with fresh labels and fresh samples, using independent labels for belief and model responses.

Feedback also closes for record likelihoods that depend only on the base and are applied to every agent. In the reciprocal potential a record is an assimilation event: it updates the laws, and the potential itself is unchanged. Every child has the same base marginal $f_0(r)$, so a record with likelihood $L_e(o\mid r)$ has the common normalization $Z_o=\int f_0L_e dA$. Its exact component update is
\begin{equation}
 f^+(r,u)=\frac{L_e(o\mid r)f(r,u)}{Z_o}.
\end{equation}
Every value at that base point receives the same multiplier, which commutes with every tree retention. The next experiment may depend on the previous record history. Induction then preserves the full finite history law, with the original continuous dynamics between updates. These admitted records see only the root marginal $f_0$. A record that reads a public sign, or a record delivered to only one subtree, multiplies the two branches by different factors and exposes the phase that the parent discarded.

Reusing one hidden child for two samples changes the response. For a pair $q_\pm=p(1\pm a\sigma)$ it gives
\begin{equation}
 p(\xi)p(\xi')[1+a(\xi)a(\xi')\sigma(\xi)\sigma(\xi')].
\end{equation}
The cross term reads the discarded public pattern. Unrestricted external probes can also expose it. Exact inheritance is a statement about the declared interface, not every experiment one could perform on the children.

The discarded contrast is also transient. If the profile is close to one that does not depend on the bits, then to first order each contrast mode decays like $e^{-\lambda t}$ with $\lambda\ge1$, and at the uniform source rows every mode decays at the common rate $3+2/(N-1)$. The same reciprocal flow that preserves the compression drives the distinctions it hides toward zero.

Finally, a phase rotation changes the public laws continuously in total variation: the angular region whose sign changes is proportional to the rotation angle. The corresponding squared Hellinger distance has the same linear dependence, so the Fisher--Rao distance grows like the square root of the angle. A phase path is therefore H\"older of order one half and has infinite Fisher--Rao length; it has no finite Fisher tangent. The original time evolution remains smooth on each fixed-partition profile chart. This construction establishes recursive boundary inheritance on a prepared family. \Cref{cf:formation} gives a different prepared family with selected membership and inherited Fisher geometry. It uses moment-dependent channels and an extensive interface; it does not establish selection of the identity-channel profile tree used here.
